\section{Spectral phase and its trace}\label{sec:phase}

We measure a unitary by the sum of its spectral angles, with each angle
chosen in $[0,1)$.  The value zero at the cut is essential: the resulting
quantity is subadditive and vanishes only at the identity.

Let $A$ be an arbitrary group, with identity $1$.  On $\ell^2(A)$ define
the left and right translations by
\[
 L(a)\delta_x=\delta_{ax},\qquad R(a)\delta_x=\delta_{xa^{-1}},
\]
and let $\mathcal M$ be the commutant of all $R(a)$.  Thus $L(a)\in\mathcal M$.
For $T\in\mathcal M$, set $\tau(T)=T_{1,1}$, and for a finite matrix
$(T_{ij})\in M_b(\mathcal M)$ set
\[
 \tau_b((T_{ij}))=\sum_{j=1}^b\tau(T_{jj}).
\]
Throughout, the matrix trace is unnormalized: $\tau_b(I)=b$.

We record why these are faithful traces, including when $A$ is uncountable.
Writing $t_x=T_{x,1}$, commutation with right translations gives
$T_{x,y}=t_{xy^{-1}}$, and $(t_x)_{x\in A}\in\ell^2(A)$.
For $S,T\in\mathcal M$,
\[
 \tau(ST)=\sum_{x\in A}s_{x^{-1}}t_x
 =\sum_{x\in A}t_{x^{-1}}s_x=\tau(TS).
\]
Both sums converge absolutely by Cauchy--Schwarz; square-summable
families have countable support even if $A$ is uncountable.  Positivity
is immediate.  If $T\geq0$ and $\tau(T)=0$, then
$T^{1/2}\delta_1=0$.  Since $T^{1/2}$ also commutes with right
translations, it kills every basis vector, so $T=0$.
Block multiplication proves traciality of $\tau_b$; the same square-root
argument, applied to each $e_j\otimes\delta_1$, proves its faithfulness.

The algebra $M_b(\mathcal M)$ is the commutant of the simultaneous right
translations on $\C^b\otimes\ell^2(A)$.  For background on group von Neumann algebras and functional calculus, see
\cite[Sections 2.2, 3.3, and 3.4, EP1]{Jones2009}. It contains spectral
projections and bounded Borel functions of its normal elements.
For a unitary on an arbitrary Hilbert space, the spectral calculus can
be obtained by decomposing the space into orthogonal cyclic reducing
subspaces: each is the closed span of $\{U^j\xi:j\in\Z\}$ and is
separable, and the calculi on these subspaces form a direct sum.
Uniqueness of the spectral resolution shows that the calculus commutes
with every unitary commuting with $U$, so it remains in the
right-translation commutant.
Moreover, $\tau_b$ is the finite sum of the vector functionals at
$e_j\otimes\delta_1$.  In particular, for a fixed unitary $U$, uniformly
bounded Borel functions $f_r$ converging pointwise to $f$ satisfy
\begin{equation}\label{eq:phase-bounded-convergence}
 \tau_b(f_r(U))\longrightarrow\tau_b(f(U)).
\end{equation}
Indeed, apply dominated convergence to the finite scalar spectral
measures associated with these vectors.  This also explains why no
separability assumption is needed.

For $\T=\{z\in\C:|z|=1\}$, define the Borel function
\[
 h(e^{2\pi ix})=x-\lfloor x\rfloor,\qquad h(1)=0,
 \qquad \ell(U)=\tau_b(h(U)).
\]
Thus $0\leq h(U)\leq I$ and $e^{2\pi i h(U)}=U$.  Faithfulness implies
\begin{equation}\label{eq:phase-faithfulness}
 \ell(U)=0\quad\Longleftrightarrow\quad U=I.
\end{equation}

\begin{lemma}[Subadditivity of spectral phase]\label{lem:phase}
For all unitaries $U,V\in M_b(\mathcal M)$,
\[
 \ell(UV)\leq\ell(U)+\ell(V).
\]
\end{lemma}

\begin{proof}
Put $B=h(V)$ and $W_u=Ue^{2\pi iuB}$ for $0\leq u\leq1$.
Then $W_0=U$, $W_1=UV$, and $W'_u=2\pi iW_uB$.
For a smooth function $f$ on $\T$, let $f'$ denote differentiation
in normalized angle:
$f'(e^{2\pi ix})=\frac{d}{dx}f(e^{2\pi ix})$.
We first claim that
\begin{equation}\label{eq:phase-derivative}
 \frac{d}{du}\tau_b(f(W_u))=\tau_b(f'(W_u)B).
\end{equation}
For $f(z)=z^r$ with $r>0$, the product rule and traciality give
\[
 \frac{d}{du}\tau_b(W_u^r)
 =2\pi i\sum_{j=0}^{r-1}\tau_b(W_u^{j+1}BW_u^{r-1-j})
 =2\pi ir\tau_b(W_u^rB).
\]
For negative powers, use $(W_u^{-1})'=-2\pi iBW_u^{-1}$ and the
same calculation.  The formula is also immediate for $r=0$.
If $f(e^{2\pi ix})=\sum_{r\in\Z}a_re^{2\pi irx}$ is smooth, then
$\sum_r(1+|r|)|a_r|<\infty$.  Its Fourier series and differentiated
series converge uniformly; the differentiated traces are uniformly
bounded termwise by $2\pi b|r a_r|\lVert B\rVert$.
Termwise differentiation proves \eqref{eq:phase-derivative}.

Suppose now that $f$ is real and $f'\leq1$.  Since $B\geq0$,
traciality and positivity give
\[
 \tau_b(f'(W_u)B)
 =\tau_b(B^{1/2}f'(W_u)B^{1/2})\leq\tau_b(B).
\]
This argument does not require $B$ to commute with $f'(W_u)$.
Integrating \eqref{eq:phase-derivative} yields
\begin{equation}\label{eq:phase-smooth-bound}
 \tau_b(f(UV))\leq\tau_b(f(U))+\tau_b(h(V)).
\end{equation}

To pass to the discontinuous function $h$, choose a smooth nonnegative
density $\rho_\varepsilon$ of integral one supported in $(0,\varepsilon)$,
where $0<\varepsilon<1$, and set
\[
 f_\varepsilon(e^{2\pi ix})
 =\int_0^\varepsilon\rho_\varepsilon(s)
       \bigl(x+s-\lfloor x+s\rfloor\bigr)\,ds.
\]
These functions are smooth and lie between zero and one.
For the periodic fractional-part function $g(x)=x-\lfloor x\rfloor$,
\[
 g(x+\delta)-g(x)
 =\delta-\bigl(\lfloor x+\delta\rfloor-\lfloor x\rfloor\bigr)
 \leq\delta\qquad(\delta\geq0).
\]
Averaging and differentiating gives $f_\varepsilon'\leq1$.
Also $f_\varepsilon\to h$ pointwise.  At the only discontinuity,
the forward shift gives
\[
 0\leq f_\varepsilon(1)
 =\int_0^\varepsilon s\rho_\varepsilon(s)\,ds
 \leq\varepsilon,
\]
so the limiting value is precisely $h(1)=0$.
Apply \eqref{eq:phase-smooth-bound} to $f_\varepsilon$ and use
\eqref{eq:phase-bounded-convergence} for the fixed unitaries $U$ and $UV$.
\end{proof}

\begin{remark}
For finite matrices, the same inequality follows from Thompson's
exponential formula, recalled in
\cite[Section 2]{AntezanaLarotondaVarela2012}. It expresses $UV$ as
$e^{2\pi i C}$ with $C$ a sum of unitary conjugates of $h(U)$ and $h(V)$.
Thus $C\geq0$, $h(e^{2\pi i C})\leq C$, and taking traces proves the bound.
The norm-limit extension to infinite factors in
\cite[Theorem 4.4]{AntezanaLarotondaVarela2012} assumes embeddability in
the ultrapower of the hyperfinite $\mathrm{II}_1$ factor. The direct
smoothing proof above requires no such approximation hypothesis and
retains the prescribed value at the phase cut.
\end{remark}
